\apendice{Estrutura e Especificação dos Artefatos Reprodutíveis}\label{ap:artefatos}

Este apêndice detalha a organização dos artefatos computacionais disponibilizados no repositório público do projeto. O módulo de avaliação (\texttt{avaliacao/}) é estruturado em \textit{runners} de experimento, conjuntos de referência versionados, variantes de \textit{prompt} e saídas por experimento --- elementos omitidos do corpo textual para preservar o registro acadêmico, mas necessários à replicação dos resultados reportados nos Capítulos~3 e~4.

\section{Correspondência nomenclatura e arquivos}

\begin{table}[htbp]
\centering
\caption{Mapeamento entre nomes usados no corpo e arquivos versionados}
\label{tab:map_artefatos}
\small
\begin{tabular}{p{6.5cm}p{6.5cm}}
\toprule
\textbf{Nome no corpo do texto} & \textbf{Arquivo versionado} \\
\midrule
Conjunto de referência SQL (80 casos) & \texttt{golden\_dataset\_v1.0.csv} \\
Conjunto de referência RAG ampliado ($n=40$) & \texttt{golden\_rag\_v1.3.csv} \\
Amostra temática curada ($n=8$) & \texttt{golden\_rag\_v1.1.csv} \\
\bottomrule
\end{tabular}
\end{table}

\section{Estrutura principal}

\begin{itemize}
    \item \texttt{avaliacao/} --- módulo de avaliação reprodutível (SQL e RAG);
    \item \texttt{avaliacao/eval/} --- \textit{runners} de experimento e biblioteca de métricas;
    \item \texttt{avaliacao/dados/golden/} --- conjunto de referência SQL (80 casos);
    \item \texttt{avaliacao/dados/golden\_rag/} --- conjunto de referência RAG ($n=40$);
    \item \texttt{avaliacao/experimentos/} --- saídas versionadas por experimento (\texttt{checkpoint.jsonl}, \texttt{metrics\_summary\_v3.json});
    \item \texttt{avaliacao/eval/prompts/} --- variantes de \textit{prompt} (compacto, estendido, regras de domínio);
    \item \texttt{scripts\_python/} --- ETL e API analítica da instância piloto;
    \item \texttt{docs/} --- documentação técnica, métricas e textos acadêmicos.
\end{itemize}

\section{Metadados do conjunto SQL}

\begin{itemize}
    \item Arquivo principal: \texttt{golden\_dataset\_v1.0.csv} (conjunto de referência SQL, 80 casos);
    \item Campo \texttt{colunas\_resposta}: grandezas avaliadas na equivalência funcional;
    \item Campo \texttt{origem}: \texttt{base\_10} (intenção analítica) ou \texttt{derivacao} (variação paramétrica);
    \item Resultados de referência: \texttt{reference\_results/NNN.json}.
\end{itemize}

\section{Metadados do conjunto RAG}

\begin{itemize}
    \item Arquivo principal: \texttt{golden\_rag\_v1.3.csv} (conjunto de referência RAG ampliado, $n=40$);
    \item Campo \texttt{doc\_id\_esperado}: rótulo canônico do ato no Diário Oficial;
    \item Campo \texttt{estrato}: \texttt{com\_numero} ou \texttt{tematico};
    \item Índice Qdrant: \textit{collection} \texttt{jornais\_caete}; correspondência por \texttt{page\_content}.
\end{itemize}

\section{Execução da avaliação}

A reprodução dos experimentos é automatizada pelo script \texttt{run\_metrics\_pipeline.ps1}, no diretório \texttt{avaliacao/}. Na forma padrão, o \textit{pipeline} executa a regressão SQL (Lote~I, experimento E1); com o parâmetro \texttt{-WithRag}, inclui também a avaliação do subsistema de recuperação documental (Lote~III). A configuração central reside em \texttt{avaliacao/eval/config.yaml}. Os requisitos de ambiente são: instância DuckDB local sobre as bases consolidadas, serviço Qdrant acessível na porta \texttt{6333} e chave de API OpenAI para \textit{embeddings} na avaliação RAG. O script completo, dependências e instruções de instalação estão versionados no repositório público do projeto.
